\documentclass[12pt]{article}
\usepackage{amsmath}
\usepackage{amssymb}
\usepackage{geometry}

\geometry{margin=1in}

\begin{document}

\begin{center}

{\fontsize{16}{18}\selectfont\textbf{CSE 400: Fundamentals of Probability in Computing}}

\vspace{10pt}

{\fontsize{14}{15}\selectfont\textbf{Lecture 20: Linear Transformations and Expectations of Multiple Random Variables}}

\vspace{8pt}
\end{center}

\hrule
\vspace{0.5cm}

\section{Outline}

\begin{itemize}
\item \textbf{Statistical Representation of Random Vectors}
\item \textbf{Mean Vector, Correlation, and Covariance Matrices}
\item \textbf{Linear Transformations}
\item \textbf{Mapping properties under $Y = AX + b$}
\item \textbf{Multivariate Gaussian Distributions}
\item \textbf{From 1D/2D cases to N-dimensional Joint PDFs}
\item \textbf{Special Cases and Properties}
\item \textbf{Implications of zero correlation and matrix factorization}
\end{itemize}

\hrule
\vspace{0.5cm}

\section{Random Vector Representation}

\subsection{Vector Form of Multiple Random Variables}

For random variables:

$X_1, X_2, ..., X_N$

We represent them as a vector:

\[
\boxed{X = [X_1, X_2, ..., X_N]^T}
\]

\subsubsection*{\textbf{Definition}}

$X$ is called a \textbf{random vector}.

\subsubsection*{\textbf{Properties}}

\begin{itemize}
\item \textbf{Provides a compact representation of multiple random variables.}
\end{itemize}

\section{Mean Vector}

\subsection{Expectation of a Random Vector}

\subsubsection*{\textbf{Mean Vector Definition}}

The mean vector of $X$ is:

\[
\boxed{\mu_X = E[X]}
\]

\subsubsection*{\textbf{Interpretation}}

\begin{itemize}
\item \textbf{Represents the average value of the random vector.}
\item \textbf{Each element corresponds to the mean of one random variable.}
\end{itemize}

\subsubsection*{\textbf{Structure}}

\[
\boxed{\mu_X = [E[X_1], E[X_2], ..., E[X_N]]^T}
\]

\section{Correlation Matrix}

\subsection{Correlation Between Multiple Random Variables}

\subsubsection*{\textbf{Definition}}

\[
\boxed{R_{XX} = E[XX^T]}
\]

Where:

\[
R_{XX}(i,j) = E[X_i X_j]
\]

\subsubsection*{\textbf{Matrix Representation}}

\[
R_{XX} =
\begin{bmatrix}
E[X_1X_1] & E[X_1X_2] & ... & E[X_1X_N] \\
E[X_2X_1] & E[X_2X_2] & ... & E[X_2X_N] \\
... & ... & ... & ... \\
E[X_NX_1] & E[X_NX_2] & ... & E[X_NX_N]
\end{bmatrix}
\]

\subsubsection*{\textbf{Property}}

\textbf{$R_{XX}$ is a symmetric matrix.}

\section{Covariance Matrix}

\subsection{Definition}

\[
\boxed{C_{XX} = E[(X - \mu_X)(X - \mu_X)^T]}
\]

\subsubsection*{\textbf{Example Form}}

\[
C_{XX} =
\begin{bmatrix}
Var(X_1) & Cov(X_1,X_2) \\
Cov(X_2,X_1) & Var(X_2)
\end{bmatrix}
\]

\subsubsection*{\textbf{Property}}

\[
Cov(X_i,X_j) = Cov(X_j,X_i)
\]

\textbf{Thus the matrix is symmetric.}

\section{Linear Transformation of Random Vectors}

\subsection{Model}

Consider the transformation:

\[
\boxed{Y = AX + b}
\]

Where:

\begin{itemize}
\item \textbf{$A$ : Transformation matrix}
\item \textbf{$b$ : Constant vector}
\item \textbf{$X$ : Input random vector}
\end{itemize}

\subsubsection*{\textbf{Interpretation}}

\textbf{$A$ scales and rotates the data.}

\textbf{$b$ shifts the center.}

\subsubsection*{\textbf{Element-wise Representation}}

\[
Y_1 = a_{11}X_1 + a_{12}X_2 + ... + a_{1N}X_N + b_1
\]

\[
Y_2 = a_{21}X_1 + a_{22}X_2 + ... + a_{2N}X_N + b_2
\]

...

\[
Y_N = a_{N1}X_1 + a_{N2}X_2 + ... + a_{NN}X_N + b_N
\]

\section{Mean of Transformed Vector}

For $Y = AX + b$:

\[
\mu_Y = E[Y]
\]

\subsection{Derivation}

Substitute $Y$:

\[
\mu_Y = E[AX + b]
\]

Apply linearity:

\[
\mu_Y = E[AX] + E[b]
\]

Since $b$ is constant:

\[
E[b] = b
\]

Thus:

\[
\boxed{\mathbf{\mu_Y = A\mu_X + b}}
\]

\section{Correlation Matrix Transformation}

Definition:

\[
R_{YY} = E[YY^T]
\]

Substitute $Y = AX + b$:

\[
R_{YY} = E[(AX + b)(AX + b)^T]
\]

\subsection{Expansion}

\[
R_{YY} = E[AXX^TA^T + AXb^T + bX^TA^T + bb^T]
\]

\subsection{Applying Expectation}

\[
R_{YY} = AE[XX^T]A^T + AE[X]b^T + bE[X]^TA^T + bb^T
\]

\subsection{Final Result}

\[
\boxed{\mathbf{R_{YY} = AR_{XX}A^T + A\mu_Xb^T + b\mu_X^TA^T + bb^T}}
\]

\section{Covariance Matrix Transformation}

Key observation:

\[
Y - \mu_Y = (AX + b) - (A\mu_X + b)
\]

\[
Y - \mu_Y = A(X - \mu_X)
\]

\subsection{Derivation}

\[
C_{YY} = E[(Y - \mu_Y)(Y - \mu_Y)^T]
\]

Substitute:

\[
C_{YY} = E[A(X - \mu_X)(X - \mu_X)^TA^T]
\]

\subsection{Result}

\[
\boxed{\mathbf{C_{YY} = AC_{XX}A^T}}
\]

\subsection{Observation}

\textbf{Shift $b$ does not affect covariance.}

\section{Summary of Linear Transformation}

If $Y = AX + b$:

Mean:

\[
\boxed{\mathbf{\mu_Y = A\mu_X + b}}
\]

Covariance:

\[
\boxed{\mathbf{C_{YY} = AC_{XX}A^T}}
\]

Correlation:

\[
\boxed{\mathbf{R_{YY} = AR_{XX}A^T + A\mu_Xb^T + b\mu_X^TA^T + bb^T}}
\]

\section{Numerical Example}

Given transformation:

\[
\boxed{Y = AX}
\]

Compute:

\begin{itemize}
\item \textbf{Mean vector of $Y$}
\item \textbf{Correlation matrix of $Y$}
\item \textbf{Covariance matrix of $Y$}
\end{itemize}

\section{Solution: Mean}

\[
\mu_Y = A\mu_X
\]

Example calculation:

\[
\mu_Y =
\begin{bmatrix}
1 & 0 & -1 \\
0 & -1 & 1 \\
-1 & 1 & 0
\end{bmatrix}
*
\begin{bmatrix}
2 \\ -1 \\ 1
\end{bmatrix}
\]

Component calculations:

1st element:

\[
(1)(2) + (0)(-1) + (-1)(1)
\]

2nd element:

\[
(0)(2) + (-1)(-1) + (1)(1)
\]

3rd element:

\[
(-1)(2) + (1)(-1) + (0)(1)
\]

Result:

\[
\boxed{\mathbf{
\mu_Y =
\begin{bmatrix}
1 \\ 2 \\ -3
\end{bmatrix}
}}
\]

\section{Correlation Matrix Solution}

Use:

\[
R_{YY} = AR_{XX}A^T
\]

Final:

\[
\boxed{\mathbf{
R_{YY} =
\begin{bmatrix}
17 & -6 & -9 \\
-6 & 14 & -8 \\
-9 & -8 & 19
\end{bmatrix}
}}
\]

\section{Covariance Matrix Solution}

Use relation:

\[
C_{YY} = R_{YY} - \mu_Y \mu_Y^T
\]

Outer Product:

\[
\mu_Y \mu_Y^T =
\begin{bmatrix}
1 & 2 & -3 \\
2 & 4 & -6 \\
-3 & -6 & 9
\end{bmatrix}
\]

Subtract:

\[
C_{YY} =
\begin{bmatrix}
17 & -6 & -9 \\
-6 & 14 & -8 \\
-9 & -8 & 19
\end{bmatrix}
-
\begin{bmatrix}
1 & 2 & -3 \\
2 & 4 & -6 \\
-3 & -6 & 9
\end{bmatrix}
\]

\[
\boxed{\mathbf{C_{YY} \text{ Final covariance matrix obtained after subtraction.}}}
\]

\section{Homework Assignment}

Given matrices $C_{XX}$ and $R_{XX}$.

Find mean vector $\mu_X$.

Relation:

\[
\boxed{C_{XX} = R_{XX} - \mu_X \mu_X^T}
\]

Rearrange:

\[
\boxed{\mu_X \mu_X^T = R_{XX} - C_{XX}}
\]

\textbf{Diagonal elements reveal components of $\mu_X$.}

\section{Gaussian Distributions}

\subsection{Univariate Gaussian}

\[
\boxed{
f_X(x) =
\left(\frac{1}{\sqrt{2\pi\sigma^2}}\right)
\exp\left( -\frac{(x - \mu)^2}{2\sigma^2} \right)
}
\]

\subsection{Joint Gaussian (2D)}

\[
\boxed{
f_{XY}(x,y) =
\frac{1}{2\pi \sigma_x \sigma_y \sqrt{(1 - \rho^2)}}
\exp\left(
-\frac{1}{2(1-\rho^2)}
\left[
\frac{(x-\mu_x)^2}{\sigma_x^2}
- \frac{2\rho(x-\mu_x)(y-\mu_y)}{\sigma_x\sigma_y}
+ \frac{(y-\mu_y)^2}{\sigma_y^2}
\right]
\right)
}
\]

\section{Multivariate Gaussian Distribution}

For random vector:

\[
X = [X_1, X_2, ..., X_N]^T
\]

\subsection{General PDF}

\[
\boxed{
f_X(x) =
\frac{1}{\sqrt{(2\pi)^N det(C_{XX})}}
\exp\left(
-\frac{1}{2}
(x - \mu_X)^T
C_{XX}^{-1}
(x - \mu_X)
\right)
}
\]

\subsection{Components}

\textbf{$\mu_X$ : Mean vector}

\textbf{$C_{XX}$ : Covariance matrix}

\section{Uncorrelated Gaussian Variables}

If:

\[
Cov(X_i, X_j) = 0 \quad \text{for} \quad i \neq j
\]

Then covariance matrix becomes diagonal.

Result:

\[
\boxed{f_X(x) = \prod f_{X_i}(x_i)}
\]

\textbf{Thus independent Gaussian components appear.}

\section{Homework Proof}

\[
\boxed{\mathbf{Var[X + Y] = Var[X] + Var[Y] + 2Cov(X,Y)}}
\]

\section{Case Study: Smart City Traffic + Weather System}

Random vector variables:

$X_1$ : Traffic flow

$X_2$ : Average speed

$X_3$ : Rainfall intensity

$X_4$ : Temperature

$X_5$ : Air pollution index

Thus:

\[
\boxed{X = [X_1, X_2, X_3, X_4, X_5]^T}
\]

\section{Mean Vector Example}

Average values example:

Traffic = 120

Speed = 35

Rainfall = 2

Temperature = 32

AQI = 180

\[
\boxed{\mu = E[X]}
\]

\section{Covariance Interpretation}

Examples:

\[
Cov(X_1,X_3) > 0
\]

\textbf{Rain increases congestion}

\[
Cov(X_2,X_3) < 0
\]

\textbf{Rain reduces speed}

\[
Cov(X_1,X_5) > 0
\]

\textbf{Traffic increases pollution}

\[
Cov(X_4,X_5) < 0
\]

\textbf{Temperature disperses pollution}

\section{Joint Gaussian Model}

Assume:

\[
\boxed{X \sim N(\mu, \Sigma)}
\]

Reasons:

\begin{itemize}
\item \textbf{Aggregation of many small effects}
\item \textbf{Sensor noise Gaussian}
\item \textbf{Closed form analytics}
\end{itemize}

System defined by:

\textbf{Mean $\rightarrow$ typical conditions}

\textbf{Covariance $\rightarrow$ interdependencies}

\section{Linear Transformation Example}

Define congestion index:

\[
\boxed{Y = 0.6X_1 - 0.4X_2 + 0.8X_3}
\]

This is linear transformation:

\[
\boxed{Y = a^T X}
\]

If:

\[
X \sim N(\mu, \Sigma)
\]

Then:

\textbf{$Y$ is Gaussian}

Mean:

\[
\boxed{E[Y] = a^T \mu}
\]

Variance:

\[
\boxed{Var(Y) = a^T \Sigma a}
\]

\section{Key Takeaways}

\begin{itemize}
\item \textbf{Random vectors represent multiple variables.}
\item \textbf{Mean and covariance describe system behavior.}
\item \textbf{Linear transformations modify mean and covariance.}
\item \textbf{Multivariate Gaussian models high-dimensional uncertainty.}
\item \textbf{Uncorrelated Gaussian variables factorize into independent PDFs.}
\end{itemize}

\hrule
\vspace{0.5cm}

\textbf{End of Lecture 20 Scribe.}

\end{document}